%!TEX root = ../main/main.tex

Regular expressions are a standard tool for locating information in
text. In extraction-oriented applications, however, recognising that a
pattern occurs is not enough: one needs the positions of the captured
substrings. Under the \emph{all-match} semantics, a query can produce
many mappings from variables to spans for the same document, including
overlapping matches and different capture assignments that start at the
same position. Systems such as REmatch implement this semantics and
must enumerate all results efficiently~\cite{RiverosJV23}.

Regular document spanners provide a formal model for these
queries~\cite{FlorenzanoRUVV20}. A spanner maps each document to a set
of mappings, and variable-set automata (VAs) recognise spanners by
reading document letters and executing variable markers. Extended VAs
(eVAs) group marker operations into capture steps, while logical-VAs
add character classes and substring search, as used by REmatch. For
deterministic sequential eVAs, enumeration can be done with
output-linear delay after preprocessing linear in the product of the
automaton and document sizes~\cite{FlorenzanoRUVV20}.

The enumeration guarantees are stated for deterministic automata.
Determinising a VA can produce exponentially many states, and the size
of the result depends on the nondeterministic automaton from which the
determinisation starts. Two classical constructions for ordinary
regular expressions are Glushkov and Thompson~\cite{GiammarresiPW98}.
Glushkov assigns a state to each symbol occurrence and uses the first,
last, and follow relations; Thompson builds the automaton inductively
and uses empty transitions. Both recognise the same language, but they
need not have the same number of reachable subsets after
determinisation. For extraction, the comparison must preserve capture
information: markers do not consume input, but they change the produced
mappings, so they cannot be eliminated as ordinary empty transitions.

This paper develops the formal basis for comparing the two
constructions in the presence of capture variables. We present the
chain VA $\to$ eVA $\to$ logical-VA with explicit conventions for runs,
validity, and induced mappings. We then define a Glushkov construction
with capture variables: the source regex formula is translated into an
ordinary expression over letters and markers, the position construction
is applied to that expression, and marker transitions are interpreted
as zero-width VA transitions. We prove that the resulting VA
$G_\gamma$ denotes exactly the same document spanner as the source
formula $\gamma$ (Theorem~\ref{theo:glushkov-correct}). The proof
separates the two-layer semantics of regex formulas from the word
semantics of the marker translation, and relates the latter to the
valid runs of $G_\gamma$.

The correctness result reduces the comparison with Thompson to a
well-defined state-count question: for a query $\gamma$, compare the
number of reachable states obtained by determinising $G_\gamma$ and the
corresponding Thompson construction $T_\gamma$, using the same
determinisation procedure, preserving capture semantics, and before
minimisation. We formulate this comparison for a class $\mathcal C$ of
expressions and discuss the choice of $\mathcal C$. Well-designed REQL
queries are the main candidate, because they guarantee sequentiality;
one-unambiguity is a natural special case, since a one-unambiguous
expression has a deterministic Glushkov
automaton~\cite{BruggemannKleinW98}.

Section~\ref{sec:preliminaries} introduces the models and the
translation $\tau$. Section~\ref{sec:results} presents the Glushkov
construction and its correctness. Section~\ref{sec:conclusions}
discusses the consequences for the comparison and the questions that
remain open.
